The conditions give $p(1)=-1$, $p(-1)=1$, and $p'(\pm1)=p''(\pm1)=0$. Thus the degree-four derivative is
\[
p'(x)=k(x-1)^2(x+1)^2=k(x^4-2x^2+1).
\]
Integration gives $p(x)=k(\dfrac{x^5}{5}-\dfrac{2x^3}{3}+x)+C$. Adding the two prescribed endpoint values gives $C=0$; at $x=1$, $\dfrac{8}{15}k=-1$, so $k=-\dfrac{15}{8}$. Therefore
\[
\boxed{p(x)=-\dfrac{3}{8}x^5+\dfrac{5}{4}x^3-\dfrac{15}{8}x}.
\]
It has the required values and vanishing first two derivatives, which are equivalent to the stated triple-root divisibilities. The derivation also proves uniqueness.
